\subsection{RAG Experimental Methodology}
\label{sec:rag_methodology}

We evaluated RAG over approximately 24,000
documents using 125 development-split queries
and supplied query--document relevance
judgements (qrels). For exact dense retrieval,
we measured Recall@$k$, Precision@$k$,
Hit@$k$, MRR@10, and NDCG@10, varying
retrieval depth over
$k\in\{1,2,3,5,10,20,50,100\}$.
These metrics assess evidence retrieval rather
than generated-answer correctness. Retrieval
timing excludes query encoding.

For generation, we compared fixed context
sizes $N\in\{1,3,5,10,20\}$, two adaptive
similarity-gap policies (thresholds 0.015
and 0.025), and an oracle condition using
gold-relevant evidence, capped at 10 records.
The oracle is a diagnostic upper-reference
condition, not a deployable retrieval method.
All conditions used the same 125 queries.
The generator was
\texttt{EuroLLM-22B-Instruct-2512}
through the ELM API.

An automated EuroLLM-based judge scored
correctness, grounding, and completeness
from 0 to 4 each; the combined score was
normalised by 12. We report mean scores,
95\% confidence intervals, paired differences,
and context-token costs. A separate
20-query, single-worker timing subset was
used for latency comparisons to reduce
concurrency-related confounding.

Each query--condition pair was evaluated
with one generation and one judgement.
A separate 40-query repeat regenerated
and rejudged answers for the five-document
condition; identical answers were not
independently rejudged. Automated judgements,
finite query coverage, and judge variability
limit the strength of small quality
differences. The earlier 30-query GPT-4.1
pilot used a different evaluation setup
and is not pooled with this study.